% GENERATED FILE. Edit canonical lessons, then run npm run generate:books.
% canonical-source-hash: fnv1a64:59269ce173f5fb37
% canonical-lessons: KA-C364-dinacari, KA-C364-atmakathe, KA-C364-jivanacaritre, KA-C364-pithike, KA-C364-tirulu

\chapter{Writing Lives}
\label{ch:ka-writing-lives}
\hlchaptermodality{364}

\begin{chapteropening}
\textbf{By the end of this chapter:} \emph{I can say a diary, an autobiography, a biography, a preface and the gist.}

Complete the last lesson of chapter 364: I can say a diary, an autobiography, a biography, a preface and the gist.
\end{chapteropening}

\section[\texorpdfstring{\kn{ದಿನಚರಿ} (dinacari)}{ದಿನಚರಿ (dinacari)}]{\kn{ದಿನಚರಿ} (dinacari) — a diary}

\label{lesson:KA-C364-dinacari}



\begin{quote}
Before the new one: say the Kannada for to emerge, then the Kannada for to become visible.
\end{quote}

\subsection*{You'll want to know: \kn{ದಿನಚರಿ}}
\textbf{\kn{ದಿನಚರಿ}} — \emph{dinacari} — \textquotedblleft{}a diary\textquotedblright{}.

\begin{quote}
\kn{ಅಜ್ಜಿ} \kn{ಪ್ರತಿದಿನ} \kn{ರಾತ್ರಿ} \kn{ದಿನಚರಿ} \kn{ಬರೆಯುತ್ತಿದ್ದರು}.
\end{quote}

\begin{quote}
\emph{ajji pratidina rātri dinacari bareyuttiddaru.}
\end{quote}

\begin{quote}
Grandmother used to write her diary every night.
\end{quote}

Plural: \kn{ದಿನಚರಿಗಳು}, \emph{dinacarigaḷu}.

It also means a daily routine: \kn{ನನ್ನ} \kn{ದಿನಚರಿ} \kn{ಬೆಳಗ್ಗೆ} \kn{ಆರು} \kn{ಗಂಟೆಗೆ} \kn{ಶುರು}.

\begin{cousinweb}
\textbf{\kn{ದಿನಚರಿ}} (\emph{dinacari}) is borrowed from Sanskrit \emph{dinacaryā}, \textquotedblleft{}daily work\textquotedblright{} (Monier-Williams), from \emph{dina} \textquotedblleft{}a day\textquotedblright{} and \emph{caryā} \textquotedblleft{}conduct\textquotedblright{} (CDIAL 4710).
\end{cousinweb}

\begin{grammarlens}[title={What you've built}]
One more word for telling what happened.
\end{grammarlens}

\subsection*{Your turn}
\begin{itemize}
\raggedright
  \item \emph{Say it:} \emph{dinacari}
  \item \emph{Say it:} \emph{dinacari}, once more
  \item \emph{Say it:} \emph{ajji pratidina rātri dinacari bareyuttiddaru.} — Grandmother used to write her diary every night.
  \item \emph{Recall:} say the Kannada for to emerge, then the Kannada for to become visible, then say \emph{dinacari} again
\end{itemize}

\begin{tcolorbox}[breakable,colback=teal!4,colframe=teal!35!black,title={Before you move on}]
What does \kn{ದಿನಚರಿ} mean? (\textquotedblleft{}A diary\textquotedblright{}.) Say it once more.
\end{tcolorbox}
\section[\texorpdfstring{\kn{ಆತ್ಮಕಥೆ} (ātmakathe)}{ಆತ್ಮಕಥೆ (ātmakathe)}]{\kn{ಆತ್ಮಕಥೆ} (ātmakathe) — an autobiography}

\label{lesson:KA-C364-atmakathe}



\begin{quote}
Before the new one: say the Kannada for to become visible, then the Kannada for a diary.
\end{quote}

\subsection*{You'll want to know: \kn{ಆತ್ಮಕಥೆ}}
\textbf{\kn{ಆತ್ಮಕಥೆ}} — \emph{ātmakathe} — \textquotedblleft{}an autobiography\textquotedblright{}.

\begin{quote}
\kn{ಆ} \kn{ಆಟಗಾರನ} \kn{ಆತ್ಮಕಥೆ} \kn{ತುಂಬಾ} \kn{ಜನಪ್ರಿಯವಾಯಿತು}.
\end{quote}

\begin{quote}
\emph{ā āṭagārana ātmakathe tumbā janapriyavāyitu.}
\end{quote}

\begin{quote}
That player's autobiography became very popular.
\end{quote}

Plural: \kn{ಆತ್ಮಕಥೆಗಳು}, \emph{ātmakathegaḷu}.

\kn{ಆತ್ಮ}, self, and \kn{ಕಥೆ}, story: the story of oneself.

\begin{cousinweb}
\textbf{\kn{ಆತ್ಮಕಥೆ}} (\emph{ātmakathe}) is borrowed from Sanskrit, a modern compound of \emph{ātman}, \textquotedblleft{}breath, self\textquotedblright{} (CDIAL 1135), and \emph{kathā}, \textquotedblleft{}a story\textquotedblright{} (CDIAL 2704).
\end{cousinweb}

\begin{grammarlens}[title={What you've built}]
One more word for telling what happened.
\end{grammarlens}

\subsection*{Your turn}
\begin{itemize}
\raggedright
  \item \emph{Say it:} \emph{ātmakathe}
  \item \emph{Say it:} \emph{ātmakathe}, once more
  \item \emph{Say it:} \emph{ā āṭagārana ātmakathe tumbā janapriyavāyitu.} — That player's autobiography became very popular.
  \item \emph{Recall:} say the Kannada for to become visible, then the Kannada for a diary, then say \emph{ātmakathe} again
\end{itemize}

\begin{tcolorbox}[breakable,colback=teal!4,colframe=teal!35!black,title={Before you move on}]
What does \kn{ಆತ್ಮಕಥೆ} mean? (\textquotedblleft{}An autobiography\textquotedblright{}.) Say it once more.
\end{tcolorbox}
\section[\texorpdfstring{\kn{ಜೀವನಚರಿತ್ರೆ} (jīvanacaritre)}{ಜೀವನಚರಿತ್ರೆ (jīvanacaritre)}]{\kn{ಜೀವನಚರಿತ್ರೆ} (jīvanacaritre) — a biography}

\label{lesson:KA-C364-jivanacaritre}



\begin{quote}
Before the new one: say the Kannada for a diary, then the Kannada for an autobiography.
\end{quote}

\subsection*{You'll want to know: \kn{ಜೀವನಚರಿತ್ರೆ}}
\textbf{\kn{ಜೀವನಚರಿತ್ರೆ}} — \emph{jīvanacaritre} — \textquotedblleft{}a biography\textquotedblright{}.

\begin{quote}
\kn{ನಾನು} \kn{ಆ} \kn{ವಿಜ್ಞಾನಿಯ} \kn{ಜೀವನಚರಿತ್ರೆ} \kn{ಓದಿದೆ}.
\end{quote}

\begin{quote}
\emph{nānu ā vijñāniya jīvanacaritre ōdide.}
\end{quote}

\begin{quote}
I read that scientist's biography.
\end{quote}

Plural: \kn{ಜೀವನಚರಿತ್ರೆಗಳು}, \emph{jīvanacaritregaḷu}.

\kn{ಜೀವನ}, life, and \kn{ಚರಿತ್ರೆ}, an account of deeds.

\begin{cousinweb}
\textbf{\kn{ಜೀವನಚರಿತ್ರೆ}} (\emph{jīvanacaritre}) is borrowed from Sanskrit, a modern compound of \emph{jīvana}, \textquotedblleft{}life\textquotedblright{} (CDIAL 5243), and \emph{caritra}, \textquotedblleft{}going, behaviour, deeds\textquotedblright{} (CDIAL 4691).
\end{cousinweb}

\begin{grammarlens}[title={What you've built}]
One more word for telling what happened.
\end{grammarlens}

\subsection*{Your turn}
\begin{itemize}
\raggedright
  \item \emph{Say it:} \emph{jīvanacaritre}
  \item \emph{Say it:} \emph{jīvanacaritre}, once more
  \item \emph{Say it:} \emph{nānu ā vijñāniya jīvanacaritre ōdide.} — I read that scientist's biography.
  \item \emph{Recall:} say the Kannada for a diary, then the Kannada for an autobiography, then say \emph{jīvanacaritre} again
\end{itemize}

\begin{tcolorbox}[breakable,colback=teal!4,colframe=teal!35!black,title={Before you move on}]
What does \kn{ಜೀವನಚರಿತ್ರೆ} mean? (\textquotedblleft{}A biography\textquotedblright{}.) Say it once more.
\end{tcolorbox}
\section[\texorpdfstring{\kn{ಪೀಠಿಕೆ} (pīṭhike)}{ಪೀಠಿಕೆ (pīṭhike)}]{\kn{ಪೀಠಿಕೆ} (pīṭhike) — a preface; an introduction}

\label{lesson:KA-C364-pithike}



\begin{quote}
Before the new one: say the Kannada for an autobiography, then the Kannada for a biography.
\end{quote}

\subsection*{You'll want to know: \kn{ಪೀಠಿಕೆ}}
\textbf{\kn{ಪೀಠಿಕೆ}} — \emph{pīṭhike} — \textquotedblleft{}a preface; an introduction\textquotedblright{}.

\begin{quote}
\kn{ಪುಸ್ತಕದ} \kn{ಪೀಠಿಕೆಯನ್ನು} \kn{ಲೇಖಕರ} \kn{ಸ್ನೇಹಿತ} \kn{ಬರೆದನು}.
\end{quote}

\begin{quote}
\emph{pustakada pīṭhikeyannu lēkhakara snēhita baredanu.}
\end{quote}

\begin{quote}
The book's preface was written by the author's friend.
\end{quote}

The opening part that comes before the story; also the background to a talk.

\begin{cousinweb}
\textbf{\kn{ಪೀಠಿಕೆ}} (\emph{pīṭhike}) is borrowed from Sanskrit \emph{pīṭhikā}, \textquotedblleft{}a stool, a base, a pedestal\textquotedblright{} (Monier-Williams): the base a book stands on.
\end{cousinweb}

\begin{grammarlens}[title={What you've built}]
One more word for telling what happened.
\end{grammarlens}

\subsection*{Your turn}
\begin{itemize}
\raggedright
  \item \emph{Say it:} \emph{pīṭhike}
  \item \emph{Say it:} \emph{pīṭhike}, once more
  \item \emph{Say it:} \emph{pustakada pīṭhikeyannu lēkhakara snēhita baredanu.} — The book's preface was written by the author's friend.
  \item \emph{Recall:} say the Kannada for an autobiography, then the Kannada for a biography, then say \emph{pīṭhike} again
\end{itemize}

\begin{tcolorbox}[breakable,colback=teal!4,colframe=teal!35!black,title={Before you move on}]
What does \kn{ಪೀಠಿಕೆ} mean? (\textquotedblleft{}A preface; an introduction\textquotedblright{}.) Say it once more.
\end{tcolorbox}
\section[\texorpdfstring{\kn{ತಿರುಳು} (tiruḷu)}{ತಿರುಳು (tiruḷu)}]{\kn{ತಿರುಳು} (tiruḷu) — the gist; the core}

\label{lesson:KA-C364-tirulu}



\begin{quote}
Before the new one: say the Kannada for a biography, then the Kannada for a preface.
\end{quote}

\subsection*{You'll want to know: \kn{ತಿರುಳು}}
\textbf{\kn{ತಿರುಳು}} — \emph{tiruḷu} — \textquotedblleft{}the gist; the core\textquotedblright{}.

\begin{quote}
\kn{ಅವನ} \kn{ಮಾತಿನ} \kn{ತಿರುಳು} \kn{ಇಷ್ಟೇ}, \kn{ಎಲ್ಲರೂ} \kn{ಸೇರಿ} \kn{ಕೆಲಸ} \kn{ಮಾಡಬೇಕು}.
\end{quote}

\begin{quote}
\emph{avana mātina tiruḷu iṣṭē, ellarū sēri kelasa māḍabēku.}
\end{quote}

\begin{quote}
The gist of his talk was this: everyone must work together.
\end{quote}

Also the kernel of a fruit or a nut.

\begin{cousinweb}
\textbf{\kn{ತಿರುಳು}} (\emph{tiruḷu}) is native Dravidian: DEDR 3252 lists Kannada \emph{tiruḷ}, \emph{tiruḷe}, \textquotedblleft{}the essential part, strength, value\textquotedblright{}.
\end{cousinweb}

\begin{grammarlens}[title={What you've built}]
One more word for telling what happened.
\end{grammarlens}

\subsection*{Your turn}
\begin{itemize}
\raggedright
  \item \emph{Say it:} \emph{tiruḷu}
  \item \emph{Say it:} \emph{tiruḷu}, once more
  \item \emph{Say it:} \emph{avana mātina tiruḷu iṣṭē, ellarū sēri kelasa māḍabēku.} — The gist of his talk was this: everyone must work together.
  \item \emph{Recall:} say the Kannada for a biography, then the Kannada for a preface, then say \emph{tiruḷu} again
\end{itemize}

\begin{tcolorbox}[breakable,colback=teal!4,colframe=teal!35!black,title={Before you move on}]
What does \kn{ತಿರುಳು} mean? (\textquotedblleft{}The gist; the core\textquotedblright{}.) Say it once more.
\end{tcolorbox}
